\section*{第 \ref{ch:b2:magmatic-differentiation} 章——岩浆分异}

\begin{solution}{exo:b2:magmatic-differentiation:1}
(a) An$_{50}$ 的液相线温度为 \qty{1450}{\celsius}。第一颗晶体位于
该温度下的固相线上，介于 An$_{66.7}$（1372）与 An$_{83.3}$（1465）之间：
$66.7 + 16.6\times(1450 - 1372)/(1465 - 1372) = 66.7 + 13.9 = \qty{80.6}{\percent}$
钙长石。(b) 在 \qty{1394}{\celsius} 时，液体位于液相线上，
液相为 An$_{33.3}$；固相线上的晶体为 $66.7 + 16.6\times22/93 =
\qty{70.6}{\percent}$。由杠杆规则，结晶部分的比例为 $(50 -
33.3)/(70.6 - 33.3) = 16.7/37.3 = 0.448$。(c) 末次液体消失于
晶体成分达到 An$_{50}$ 时，温度为固相线温度 \qty{1287}{\celsius}；
此处液相线介于 An$_{11.1}$（1265）与 An$_{20}$（1334）之间：$11.1 +
8.9\times22/69 = \qty{13.9}{\percent}$ 钙长石。
\end{solution}

\begin{solution}{exo:b2:magmatic-differentiation:2}
$\bar D_{\ce{Sr}} = 0.15\times0.000063 + 0.35\times0.157 + 0.50\times2.7 = 0.00001 +
0.0550 + 1.35 = 1.40$; $\bar D_{\ce{Zr}} = 0.0001 + 0.0683 + 0.0005 = 0.069$;
$\bar D_{\ce{Rb}} = 0.00001 + 0.0012 + 0.0125 = 0.014$。只有被斜长石容纳的锶
为\omterm{def:b2:elements-in-the-earth:partition}{相容元素}（$\bar D > 1$）；锆和铷为\omterm{def:b2:elements-in-the-earth:partition}{不相容元素}，
其中铷的不相容性更强。
\end{solution}

\begin{solution}{exo:b2:magmatic-differentiation:3}
(a) $C_L/C_0 = 0.5^{D-1}$：Sr $0.5^{0.40} = 0.758$（含量下降）；Zr $0.5^{-0.931} =
1.907$；Rb $0.5^{-0.986} = 1.981$，接近极限值 2。(b) Zr/Rb 的比值变化为
$0.5^{0.069 - 0.014} = 0.5^{0.055} = 0.963$，不足 \qty{4}{\percent}，而两元素含量各自几乎翻倍：
正如推论所言，两个\omterm{def:b2:elements-in-the-earth:partition}{不相容元素}的比值保持母\omterm{def:b2:magma-emplacement:magma}{岩浆}的值，
因而反映源区，而非
结晶作用。
\end{solution}

\begin{solution}{exo:b2:magmatic-differentiation:4}
(a) \ce{SiO2} $+5.03$，\ce{MgO} $-1.95$，\ce{K2O} $+0.82$（以百分数计），氧化钾
增为原来的 $1.67/0.85 = 1.96$ 倍。(b) 对于完全\omterm{def:b2:elements-in-the-earth:partition}{不相容元素}
$C_L/C_0 = 1/F$，故 $F = 0.85/1.67 = 0.509$：约一半\omterm{def:b2:magma-emplacement:magma}{岩浆}须已
结晶。(c) 每 \qty{100}{g} 移去 \qty{12.5}{g}，剩余 $F \approx 0.88$，
则氧化钾仅升高约 $1/0.88 = 1.14$ 倍：单纯结晶作用无法
解释这一升高。氧化钾来自外部，源自花岗质捕掳体，其富含
氧化钾（\qty{3.6}{\percent}），这些捕掳体被\omterm{def:b2:magma-emplacement:magma}{岩浆}同化。
\end{solution}

\begin{solution}{exo:b2:magmatic-differentiation:5}
(a) $F = 0.8$，$D = 5$。\omterm{def:b2:magmatic-differentiation:crystallisation}{分离结晶}：$0.8^{4} = 0.410$；\omterm{def:b2:magmatic-differentiation:crystallisation}{平衡结晶}：
$1/(5 + 0.8\times(1 - 5)) = 1/1.8 = 0.556$。(b) $F = 0.1$，$D = 0.01$。\omterm{def:b2:magmatic-differentiation:crystallisation}{分离结晶}：
$0.1^{-0.99} = 9.77$；\omterm{def:b2:magmatic-differentiation:crystallisation}{平衡结晶}：$1/(0.01 + 0.1\times0.99) = 1/0.109 = 9.17$。
(c) 当 $F \to 0$ 时，瑞利定律趋于 $0$（$D > 1$ 的情形），
对 $D < 1$ 则趋于无穷，而\omterm{def:b2:magmatic-differentiation:crystallisation}{平衡结晶}定律趋于 $1/D$（末次液体与
整体成分的晶体处于平衡）。只有\omterm{def:b2:magmatic-differentiation:crystallisation}{分离结晶}才能彻底
耗尽液体中的\omterm{def:b2:elements-in-the-earth:partition}{相容元素}。
\end{solution}

\begin{solution}{exo:b2:magmatic-differentiation:6}
(a) 一个化学式单位含 1.6 \ce{MgO}、0.4 \ce{FeO} 和一个 \ce{SiO2}：$1.6\times
40.30 = 64.48$，$0.4\times71.84 = 28.74$，60.08，合计 \qty{153.30}{g}。质量
百分比：\ce{MgO} \qty{42.06}{\percent}，\ce{FeO} \qty{18.75}{\percent}，
\ce{SiO2} \qty{39.19}{\percent}。(b) 对 \ce{MgO} 应用杠杆规则：$x = (10.0 - 8.0)/(42.06
- 8.0) = 2.0/34.06 = 0.0587$，约占 \qty{6}{\percent} 的橄榄石。(c) 液体的
二氧化硅：$(49.0 - 0.0587\times39.19)/(1 - 0.0587) = 46.70/0.9413 =
\qty{49.61}{\percent}$。移去橄榄石（其二氧化硅含量低于玄武岩）会使
液体的二氧化硅含量升高，但升幅很慢。
\end{solution}

\begin{solution}{exo:b2:magmatic-differentiation:7}
(a) $F(0) = 0.768$；$F(1000) = 0.1091 - 0.5906 + 0.7678 = 0.286$；$F(2000) = 0.4364
- 1.1813 + 0.7678 = 0.023$。(b) 判别式为 $(5.9064\times10^{-4})^2 -
4\times1.091\times10^{-7}\times0.7678 = 3.4885\times10^{-7} - 3.3507\times10^{-7} =
1.379\times10^{-8}$，其平方根为 $1.174\times10^{-4}$；较小根为
$H = (5.9064 - 1.174)\times10^{-4}/(2.182\times10^{-7}) = \qty{2169}{m}$，即
三明治层。(c) $1 - 0.768 = 0.232$：接近四分之一的\omterm{def:b2:magma-emplacement:magma}{岩浆}
已在出露的下部层带底部形成之前结晶，结晶部位为其下方的隐藏带
以及沿侧壁与顶板分布的边缘系列。
\end{solution}

\begin{solution}{exo:b2:magmatic-differentiation:8}
\omterm{def:b2:magmatic-differentiation:cumulate}{堆晶岩}中含有 $D\,C_L = D\,C_0F^{D-1}$，故两个层位之间
$\ln(\mathrm{Ni}_2/\mathrm{Ni}_1) = (D - 1)\ln(F_2/F_1)$。(a) $D - 1 = \ln(54/177)/
\ln(0.29/0.70) = (-1.187)/(-0.881) = 1.35$，故 $D \approx 2.35$（
对 $F \ge 0.23$ 的 66 个\omterm{def:b2:magmatic-differentiation:cumulate}{堆晶岩}作回归，得 2.38）。(b) $D - 1 = \ln(6.2/54)/\ln(0.17/0.29) = (-2.164)/(-0.534)
= 4.05$，故 $D \approx 5.05$。(c) 总分配系数 $D$ 在结晶矿物
组合改变时发生变化（新增吸收镍的矿物，或某矿物退出），并且
\omterm{def:b2:elements-in-the-earth:partition}{矿物–熔体分配系数}本身也会随液体冷却、成分改变而升高，
如橄榄石的实测所示；\omterm{def:b2:magmatic-differentiation:cumulate}{堆晶岩}中残留的液体会使
估计变得模糊。
\end{solution}

\begin{solution}{exo:b2:magmatic-differentiation:9}
(a) $100 + 25.4 - 2.9 - 9.6 = \qty{112.9}{g}$。(b) \ce{SiO2} 的质量（克）：$55.53 +
25.4\times0.7248 - 2.9\times0.3919 - 9.6\times0.5054 = 55.53 + 18.41 - 1.14 - 4.85 =
\qty{67.95}{g}$，故 $67.95/112.9 = \qty{60.19}{\percent}$，对比 60.16。\ce{MgO}：
$4.66 + 0.20 - 1.22 = \qty{3.64}{g}$，\qty{3.22}{\percent}，对比 3.19。\ce{K2O}：
$0.88 + 0.92 = \qty{1.80}{g}$，\qty{1.60}{\percent}，对比 1.75。(c) 二氧化硅与氧化镁，为三个约束中的两个，
拟合在 0.03 以内；氧化钾未作约束，偏低 0.15，接近
其数值的十分之一：同化物质含有比平均
捕掳体更多的氧化钾，或其富钾熔体被优先带入。
\end{solution}

\begin{solution}{exo:b2:magmatic-differentiation:10}
(a) $z = (0.2 + 1.5 - 1)/(0.2 - 1) = 0.7/(-0.8) = -0.875$；$F^{-z} = 0.7^{0.875} =
0.7319$；$C^*/C_0 = 0.2\times0.5/(0.2 - 1 + 1.5) = 0.1/0.7 = 0.1429$，故 $C/C_0 =
0.7319 + 0.1429\times(1 - 0.7319) = 0.7319 + 0.0383 = 0.770$。(b) 按瑞利定律：
$0.7^{0.5} = 0.837$。AFC \omterm{def:b2:magma-emplacement:magma}{岩浆}的锶含量较低，原因有二：要使
其质量减少 \qty{30}{\percent}，同时获得 $r$ 倍于结晶量的物质，则须
结晶 $(1 - 0.7)/(1 - 0.2) = 0.375$，而非 0.30；且其获得的锶
来自富集程度仅为一半的岩石。(c) $(0.7319\times0.7035 + 0.0383\times0.7200)/0.7702
= (0.5149 + 0.0276)/0.7702 = 0.70432$：比值升高了 0.0008，而
仅结晶作用则会使其维持在 0.7035。
\end{solution}

\begin{solution}{exo:b2:magmatic-differentiation:11}
(a) $1.1\times(1100 - 200) + 250 = 990 + 250 = \qty{1240}{J}$，每克。(b)
$25.4\times1240 = \qty{31496}{J}$，接近威尔科克斯给出的 \qty{31350}{J}。(c)
$31496/370 = \qty{85.1}{g}$ 斜长石。(d) 质量平衡仅移去 $2.9 +
9.6 = \qty{12.5}{g}$ 的晶体，其潜热（约 \qty{4600}{J}）只能覆盖
所需热量的不足六分之一；$r = 25.4/12.5 = 2.03$，同化岩石的量为
所形成晶体的两倍，而热量收支只允许 $r$ 远小于 1。喷出的\omterm{def:b2:magma-emplacement:magma}{岩浆}
本身无法提供热量：必须有质量更大的\omterm{def:b2:magma-emplacement:magma}{岩浆}已经结晶，
并在深部冷却，不为人所见。
\end{solution}

\begin{solution}{exo:b2:magmatic-differentiation:12}
(a) 第 1 阶段，$F$ 从 1 降到 0.8，$D = 5$，因子为 $0.8^{4} = 0.4096$。第 2 阶段，
从 0.8 降到 0.5，$D = 1.5$：定律适用于剩余液体，因此
因子为 $(0.5/0.8)^{0.5} = 0.7906$。镍：$250\times0.4096\times0.7906 =
\qty{81.0}{ppm}$。(b) $250\times0.5^{4} = \qty{15.6}{ppm}$，低了五倍。(c)
\omterm{thm:b2:magmatic-differentiation:rayleigh}{瑞利分馏定律}假定 $D$ 为常数；当结晶出的矿物组合改变时，
$D$ 发生跳变，定律须从变化时所达到的液相成分起重新应用，
并以该液相为基准度量 $F$。
\end{solution}

\begin{solution}{pb:b2:magmatic-differentiation:1}
\noindent\textbf{第一部分。}

\textbf{1.} $\bar D_{\ce{Ni}} = 0.20\times5.0 + 0.50\times0 + 0.30\times2.0 = 1.60$。

\textbf{2.} $\bar D_{\ce{Zr}} = 0.20\times0.0007 + 0.50\times0.001 + 0.30\times0.195 =
0.00014 + 0.0005 + 0.0585 = 0.059$.

\textbf{3.} 上部带：$\bar D_{\ce{Ni}} = 0.45\times2.0 = 0.90$；$\bar D_{\ce{Zr}} =
0.55\times0.001 + 0.45\times0.195 = 0.088$.

\textbf{4.} 镍在下部带中为\omterm{def:b2:elements-in-the-earth:partition}{相容元素}（1.60），并
在橄榄石消失后略呈不相容（0.90）；锆在两个带中均为\omterm{def:b2:elements-in-the-earth:partition}{不相容元素}。

\textbf{5.} 单斜辉石：它提供了 $0.0585$，占下部带 $0.059$ 的
几乎全部，因为它是唯一分配系数远离零的矿物。

\textbf{6.} $(C_{\ce{Ni}}/C_{\ce{Zr}})$ 的变化因子为 $F^{D_{\ce{Ni}} - D_{\ce{Zr}}}$；
每减半一次，$0.5^{1.60 - 0.059} = 0.5^{1.541} = 0.344$。

\noindent\textbf{第二部分。}

\textbf{7.} $1.60\times150 = \qty{240}{ppm}$。

\textbf{8.} 液相中含量为 $190/1.60 = \qty{118.75}{ppm}$；$F = (118.75/150)^{1/0.6} =
0.7917^{1.667} = 0.68$.

\textbf{9.} $80\times0.68^{-0.941} = 80\times1.438 = \qty{115}{ppm}$。

\textbf{10.} \omterm{def:b2:magmatic-differentiation:crystallisation}{分离结晶}：$150\times0.68^{0.6} = 150\times0.793 = \qty{119.0}{ppm}$；
\omterm{def:b2:magmatic-differentiation:crystallisation}{平衡结晶}：$150/(1.6 + 0.68\times(1 - 1.6)) = 150\times0.839 = \qty{125.8}{ppm}$。
\omterm{def:b2:magma-emplacement:magma}{岩浆}结晶掉三分之一后，二者差异仅为 \qty{6}{\percent}，
对中等相容的元素而言，两种结晶方式在早期彼此接近。

\textbf{11.} $150\times0.55^{0.6} = \qty{104.8}{ppm}$。

\textbf{12.} 层位下方：$1.60\times104.8 = \qty{167.7}{ppm}$；层位上方：
$0.90\times104.8 = \qty{94.3}{ppm}$。\omterm{def:b2:magmatic-differentiation:cumulate}{堆晶岩}的镍剖面在
橄榄石自岩层中消失的层位处，含量突然下降 \qty{44}{\percent}，
呈阶跃状，而非平滑曲线。

\textbf{13.} 液相：$104.8\times(0.40/0.55)^{0.90 - 1} = 104.8\times1.032 =
\qty{108.2}{ppm}$；\omterm{def:b2:magmatic-differentiation:cumulate}{堆晶岩} $0.90\times108.2 = \qty{97.4}{ppm}$。当 $\bar D < 1$
时，镍在液相和\omterm{def:b2:magmatic-differentiation:cumulate}{堆晶岩}中缓慢升高：剖面在
阶跃之后平缓上升。

\noindent\textbf{第三部分。}

\textbf{14.} $x_{\mathrm{ol}}(m_{\mathrm{ol},j} - d_j) + x_{\mathrm{pl}}
(m_{\mathrm{pl},j} - d_j) + x_{\mathrm{cpx}}(m_{\mathrm{cpx},j} - d_j) = 100\,(p_j -
d_j)$.

\textbf{15.} \ce{MgO}：$37.16\,x_{\mathrm{ol}} - 4.9\,x_{\mathrm{pl}} +
11.1\,x_{\mathrm{cpx}} = 410$；\ce{CaO}：$-7.3\,x_{\mathrm{ol}} +
7.06\,x_{\mathrm{pl}} + 13.2\,x_{\mathrm{cpx}} = 420$；\ce{Al2O3}：
$-14.1\,x_{\mathrm{ol}} + 17.6\,x_{\mathrm{pl}} - 11.1\,x_{\mathrm{cpx}} = 190$.

\textbf{16.} 用消元法或克拉默法则求解：$x_{\mathrm{ol}} = 9.03$，
$x_{\mathrm{pl}} = 30.84$，$x_{\mathrm{cpx}} = 20.32$，单位为每 \qty{100}{g} 中的克数。

\textbf{17.} 晶体 $9.03 + 30.84 + 20.32 = \qty{60.19}{g}$；液相
\qty{39.81}{g}，$F = 0.398$。

\textbf{18.} $(4950 - 9.03\times39.19 - 30.84\times50.54 - 20.32\times51.5)/39.81 =
(4950 - 353.9 - 1558.7 - 1046.5)/39.81 = 1990.9/39.81 = \qty{50.0}{\percent}$，与
\omterm{def:b2:magma-emplacement:bodies}{岩墙}中的结果一致：未用作对照的氧化物被复现。

\textbf{19.} 被移去的晶体中，橄榄石占 $9.03/60.19 = \qty{15}{\percent}$，
斜长石占 \qty{51}{\percent}，单斜辉石占 \qty{34}{\percent}，介于
两个带的矿物组成之间。它们与从 $F = 1$ 降至
0.55 结晶出的下部组合相符（$0.45\times20 = 9.0$，$0.45\times50 = 22.5$，$0.45\times30 = 13.5$），随后是
从 0.55 降至 0.40 结晶出的上部组合（$0.15\times55 = 8.25$，$0.15\times45 = 6.75$）：9.0，
30.75 和 20.25 克。

\textbf{20.} $0.30/0.75 = 0.40$。

\noindent\textbf{第四部分。}

\textbf{21.} $0.55^{0.059 - 1} = 0.55^{-0.941} = 1.755$；$80\times1.755 =
\qty{140.4}{ppm}$。

\textbf{22.} $188/140.4 = 1.339 = (F/0.55)^{0.088 - 1}$，故 $F = 0.55\times1.339^{-1/
0.912} = 0.55\times0.726 = 0.399$.

\textbf{23.} $(188/80)^{1/(0.07 - 1)} = 2.35^{-1.075} = 0.399$。对于\omterm{def:b2:elements-in-the-earth:partition}{不相容元素}，
$C_L/C_0 \approx 1/F$，无论 $D$ 多小，
答案几乎不取决于 $D$ 的细节，因此\omterm{def:b2:elements-in-the-earth:partition}{不相容元素}是估计 $F$ 的最佳标尺。

\textbf{24.} 斜长石不含铁，占晶体的一半，而
辉石和橄榄石合计所含铁少于其所含氧化镁：液相
铁含量升高、镁含量降低，而二氧化硅并未增加。这种
在二氧化硅含量近乎恒定时的铁富集，正是
下一章所讲拉斑玄武岩系列的标志。

\textbf{25.} 纯堆晶矿物将含有 $0.088\times188 = \qty{16.5}{ppm}$。
设捕获液体的质量分数为 $t$，则 $16.5(1 - t) + 188\,t = 47$，故 $t = (47 -
16.5)/(188 - 16.5) = 30.5/171.5 = 0.18$，即约 \qty{18}{\percent} 的捕获液体。

\textbf{26.} 三个估计值，即主量元素（0.398）、钾（0.40）和
锆（0.399）三者一致，表明最上层结晶时，
原始\omterm{def:b2:magma-emplacement:magma}{岩浆}中仍呈液态的部分占
\textbf{\qty[text-series-to-math=true]{40}{\percent}}。
\end{solution}
